% =====================================================================
%  Seance 3 -- B. cours 2 : classification sur classe rare (frames 10--18)
% =====================================================================
\section{Classification supervisée sur classe rare}

\begin{frame}{Le problème : 0,35 \% de positifs}
\begin{columns}[t]
\begin{column}{0.5\textwidth}
\begin{center}
\includegraphics[width=\textwidth]{f18_desequilibre.png}
\end{center}
\end{column}
\begin{column}{0.46\textwidth}
\begin{block}{Le split temporel, seul valable}
\begin{center}\small
\begin{tabular}{@{}lll@{}}
\toprule
 & fenêtres & positifs\\
\midrule
train $<$ 22/09 & 24 192 & 56\\
test $\ge$ 22/09 & 10 368 & 64\\
\bottomrule
\end{tabular}
\end{center}
Le test contient l'incident du 22/09 à 3 h : le modèle est jugé sur
un \textbf{futur jamais vu}.
\end{block}
\begin{wbox}[L'exactitude est trompeuse]
Un modèle « toujours 0 » atteint \textbf{99,65 \%} d'exactitude
(34 440 / 34 560) --- et ne détecte \textbf{rien}. À cette rareté,
l'exactitude ne mesure plus rien : il faut la matrice de confusion.
\end{wbox}
\end{column}
\end{columns}
\end{frame}

\begin{frame}[fragile]{Les 4 modèles, entraînés pour de vrai}
\begin{codebox}[Le même protocole pour les quatre --- split temporel strict]
\begin{lstlisting}
from sklearn.linear_model import LogisticRegression
from sklearn.ensemble import (RandomForestClassifier,
                              GradientBoostingClassifier)
from sklearn.preprocessing import StandardScaler

Xtr, ytr = X[train], f.loc[train, "incident_30min"]
Xte, yte = X[~train], f.loc[~train, "incident_30min"]

rf = RandomForestClassifier(300, class_weight="balanced",
                            random_state=5534, n_jobs=2)
gb = GradientBoostingClassifier(random_state=5534)
sc = StandardScaler().fit(Xtr)          # statistiques du TRAIN
lr = LogisticRegression(class_weight="balanced",
                        max_iter=3000).fit(sc.transform(Xtr), ytr)
\end{lstlisting}
\end{codebox}
\begin{wbox}[Où la fuite se cache quand on standardise]
\py{StandardScaler().fit(X)} sur train+test mélangés importe la
moyenne \emph{du futur} dans le passé. Toujours \textbf{fit sur
le train}, \py{transform} sur les deux. Même règle pour tout
statistique appris (seuil, quantile, moyenne mobile).
\end{wbox}
\end{frame}

\begin{frame}{La matrice de confusion : le seul regard honnête}
\begin{columns}[t]
\begin{column}{0.42\textwidth}
\begin{center}\small
\begin{tabular}{@{}lp{2.0cm}p{2.0cm}@{}}
\toprule
 & prédit 0 & prédit 1\\
\midrule
réel 0 & \textbf{TN} juste & \textbf{FP} fausse alarme\\
réel 1 & \textbf{FN} manqué & \textbf{TP} détecté\\
\bottomrule
\end{tabular}
\end{center}
\vspace{1mm}
\begin{block}{Les deux métriques qui comptent}
\begin{itemize}
\item \textbf{Précision} $= TP/(TP{+}FP)$ : quand ça alerte, à quel
point c'est vrai ? $\rightarrow$ fatigue du NOC.
\item \textbf{Rappel} $= TP/(TP{+}FN)$ : quelle part des incidents
est vue ? $\rightarrow$ SLA.
\item $F_1 = 2\,\dfrac{p \times r}{p + r}$ : compromis unique.
\end{itemize}
\end{block}
\end{column}
\begin{column}{0.54\textwidth}
\begin{center}
\includegraphics[width=\textwidth]{f19_matrice_confusion.png}
\end{center}
\begin{kbox}[Mesuré sur le test]
Gradient boosting, seuil 0,50 : \textbf{40 TP, 0 FP, 24 FN} ---
précision 1,00, rappel 0,625, $F_1$ 0,769 (chiffres réels du JSON).
\end{kbox}
\end{column}
\end{columns}
\end{frame}

\begin{frame}{Les 4 familles de modèles, en une slide chacune}
\begin{columns}[t]
\begin{column}{0.48\textwidth}
\begin{block}{Régression logistique}
\begin{itemize}
\item $p(\text{incident}) = \sigma(w \cdot x)$ : linéaire, calibrée,
\textbf{interprétable} (chaque poids = effet d'une variable).
\item Rapide, stable, mais rate les frontières non linéaires.
\end{itemize}
\end{block}
\begin{block}{Arbre de décision}
\begin{itemize}
\item Règles \emph{ lisibles} : « si \py{lat\_p99} $>$ 20 et
\py{perte\_moy} $>$ 0,5 alors alerte ».
\item Seul, il sur-apprend : profond = bruit mémorisé.
\end{itemize}
\end{block}
\end{column}
\begin{column}{0.48\textwidth}
\begin{block}{Forêt aléatoire}
\begin{itemize}
\item Moyenne de 300 arbres sur échantillons/variables aléatoires :
\textbf{robuste}, peu de réglage, donne des \textbf{importances}.
\item \py{class\_weight='balanced'} : les 56 exemples du train
pèsent comme $\approx$ 12 000 négatifs.
\end{itemize}
\end{block}
\begin{block}{Gradient boosting}
\begin{itemize}
\item Arbres \emph{successifs} qui corrigent chacun les erreurs des
précédents : souvent le \textbf{meilleur} sur tabulaire.
\item Plus sensible au réglage, moins parallélisable.
\end{itemize}
\end{block}
\end{column}
\end{columns}
\end{frame}

\begin{frame}{Classe rare : \py{class\_weight='balanced'} vs SMOTE}
\begin{columns}[t]
\begin{column}{0.5\textwidth}
\begin{block}{class\_weight='balanced' (retenu)}
\begin{itemize}
\item Pèse la fonction de coût : une erreur sur la classe rare
compte $\approx \dfrac{N}{2 \times N_{min}}$ fois plus.
\item Aucune donnée synthétique : \textbf{pas de risque de fuite}
ni d'artefacts.
\item Effet mesuré ici : le seuil de décision effectif se déplace
--- à 0,50 la forêt balanced est \emph{trop} agressive
(précision 0,00 sur le test).
\end{itemize}
\end{block}
\end{column}
\begin{column}{0.46\textwidth}
\begin{block}{SMOTE (mention)}
\begin{itemize}
\item Crée des exemples minoritaires \textbf{synthétiques} par
interpolation entre voisins rares.
\item DANGER en séries temporelles : interpoler entre deux
fenêtres de \emph{serveurs différents} ou de \emph{jours différents}
fabrique des points qui ne peuvent pas exister.
\item Retenu : \py{class\_weight}, plus simple et sûr ici.
\end{itemize}
\end{block}
\end{column}
\end{columns}
\end{frame}

\begin{frame}{Quiz 2 --- classe rare (5 min)}
\begin{qbox}[3 questions]
\begin{enumerate}
\item Un modèle « toujours 0 » sur notre dataset a une exactitude de :\\
(a) 50 \% \quad (b) 99,65 \% \quad (c) 0,35 \%
\item Précision 1,0 et rappel 0,1 signifie :\\
(a) le modèle est parfait \quad (b) chaque alerte est vraie, mais
9 incidents sur 10 sont manqués \quad (c) il n'alerte jamais
\item En série temporelle, SMOTE est risqué car :\\
(a) il est lent \quad (b) il fabrique des fenêtres synthétiques
sans réalité physique ni cohérence temporelle \quad (c) il exige
des étiquettes
\end{enumerate}
\end{qbox}
\vspace{1mm}
\begin{solbox}[Réponses]
\textbf{1 : b} (34 440 / 34 560) \quad \textbf{2 : b} \quad \textbf{3 : b}
\end{solbox}
\end{frame}

\begin{frame}{Précision-rappel bat la courbe ROC ici}
\begin{columns}[t]
\begin{column}{0.48\textwidth}
\begin{center}
\includegraphics[width=\textwidth]{f20_courbes_pr.png}
\end{center}
\end{column}
\begin{column}{0.48\textwidth}
\begin{block}{Pourquoi pas la ROC / AUC ?}
La ROC utilise \textbf{FPR} $= FP/(FP{+}TN)$ : avec 10 304 négatifs,
même 300 fausses alarmes donnent un FPR minuscule $\Rightarrow$
la ROC \textbf{flatte} les modèles sur classes très rares.
\end{block}
\begin{kbox}[La courbe PR, elle, parle métier]
Chaque point $=$ un couple (fatigue NOC, SLA). La ligne pointillée
est le \textbf{hasard} (0,35 \%). Mesuré : gradient boosting
\textbf{AP = 0,63}, forêt 0,19, logistique 0,32 --- l'écart se lit
directement en « alertes utiles par jour ».
\end{kbox}
\begin{wbox}[Coût asymétrique]
\textbf{FP} : opérationnel dérangé pour rien --- fatigue,
désensibilisation. \textbf{FN} : SLA cassé (99,80 \% mesuré).
Ici le FN coûte plus cher, mais un flux de FP détruit la
confiance en l'outil.
\end{wbox}
\end{column}
\end{columns}
\end{frame}
